\documentclass[11pt,reqno]{amsart}
\usepackage[T1]{fontenc}
\usepackage{lmodern,amsmath,amscd,amssymb,amsthm,mathtools,mathrsfs,microtype,booktabs,longtable}
\usepackage[margin=1.08in]{geometry}
\usepackage{xr-hyper}
\usepackage[hidelinks]{hyperref}
\numberwithin{equation}{section}
\newtheorem{theorem}{Theorem}[section]
\newtheorem{proposition}[theorem]{Proposition}
\newtheorem{lemma}[theorem]{Lemma}
\newtheorem{corollary}[theorem]{Corollary}
\theoremstyle{definition}\newtheorem{definition}[theorem]{Definition}
\newtheorem{example}[theorem]{Example}
\theoremstyle{remark}\newtheorem{remark}[theorem]{Remark}
\newcommand{\C}{\mathbb C}\newcommand{\Q}{\mathbb Q}
\newcommand{\Pj}{\mathbb P}\newcommand{\A}{\mathbb A}
\newcommand{\OO}{\mathcal O}\newcommand{\II}{\mathcal I}
\newcommand{\NN}{\mathcal N}\newcommand{\LL}{\mathcal L}
\newcommand{\ZZ}{\mathcal Z}\newcommand{\RR}{\mathcal R}
\DeclareMathOperator{\Sym}{Sym}\DeclareMathOperator{\Gr}{Gr}
\DeclareMathOperator{\Hom}{Hom}\DeclareMathOperator{\Tot}{Tot}
\DeclareMathOperator{\Fitt}{Fitt}\DeclareMathOperator{\Ann}{Ann}
\DeclareMathOperator{\rank}{rank}\DeclareMathOperator{\Proj}{Proj}
\DeclareMathOperator{\Spec}{Spec}\DeclareMathOperator{\PGL}{PGL}
\DeclareMathOperator{\Pic}{Pic}\DeclareMathOperator{\diag}{diag}
\DeclareMathOperator{\Tor}{Tor}\DeclareMathOperator{\im}{im}
\DeclareMathOperator{\coker}{coker}\DeclareMathOperator{\codim}{codim}
\DeclareMathOperator{\End}{End}
\DeclareMathOperator{\disc}{disc}\DeclareMathOperator{\Disc}{Disc}
\DeclareMathOperator{\Supp}{Supp}\DeclareMathOperator{\Hilb}{Hilb}
\DeclareMathOperator{\rad}{rad}\DeclareMathOperator{\Ass}{Ass}
\DeclareMathOperator{\gr}{gr}\DeclareMathOperator{\Gal}{Gal}
\DeclareMathOperator{\length}{length}
\allowdisplaybreaks[2]
\setlength{\emergencystretch}{3em}

\begin{document}
\title[Finite failure schemes and quadratic pencils]{Finite failure schemes and the reconstruction of quadratic pencils}
\author{Qian Qi}
\date{September 26, 2026}
\subjclass[2020]{14M12, 14B05, 14D20, 13C40}
\keywords{Quadratic pencil, finite failure algebra, intrinsic reconstruction, common divisor, conductor, normalization fibre}
\hypersetup{pdftitle={Finite failure schemes and the reconstruction of quadratic pencils, revision 170},pdfauthor={Qian Qi}}
\begin{abstract}
We reconstruct quadratic pencils from their sharp unmarked finite
failure algebras and analyze the associated multiplication graphs.
The complete first nonreduced Hilbert fibre has reduced components
$\Pj^4$ and $\mathbb F_2$ meeting along doubled lines, together with
a square-zero nilpotent ideal. Products give all fibres with smaller
corank-two Smith exponents at most two. An explicit Hilbert--Burch
chart, a fixed-target embedded comparison, and a connectedness
argument establish the full scheme statement. This master contains
the complete proofs of both companion papers and their supplementary
results. We also determine the collision of
two reduced incidences, separating its flat horizontal closure from
its vertical conic excess and identifying the global nilradical sheaf.
Universal horizontal models, all-arc double-contact presentations,
the nonsplit algebra extension, and arbitrary ramification are
proved in the additional sections.
The determinantal model for every polynomial contact order, its finite
jet and decorated chamber theorems, and the complete nonsymmetric
two-wall slice are included without removing predecessor proofs.
The all-order marked-contact chain, its normal Rees algebra,
punctual genus-correction rings, and fixed-degree equivariant gluing
are also included in full.
The normalization, full parameter fibre, global conductor, and root-label
identifications of every marked coefficient cover are included as well.
\end{abstract}
\maketitle
\section*{The paired manuscripts}
\label{sec:introduction-v153}
The sharp inverse is the main theorem of the first paper. The second
paper's new main results are
Theorems~\ref{thm:complete-fibre-v163} and
\ref{thm:product-fibres-v163}, based on
Proposition~\ref{prop:embedded-functors-v163} and
Lemma~\ref{lem:HB-family-v163}. This unified source is a preservation
master, not a third submission. Prior proof blocks are retained,
while the focused companion sources provide the primary reading order.
The effective-family application uses reconstruction; it does not
assert an internal boundary operation before it. The Ballico 1993
full-text comparison and an external independent audit of the whole
inverse proof remain uncompleted.

\subsection{Collision families}
Theorem~\ref{thm:collision-family-v164} identifies a full discriminant
pullback and its horizontal closure. Theorem~\ref{thm:ordered-resolution-v164}
gives its ordered normalization after a quadratic base change.
Theorem~\ref{thm:nilradical-sheaf-v164} identifies the global square-zero
sheaf, and Proposition~\ref{prop:genuine-collision-v164} supplies a
fixed-target pencil realization. The parameter-space semistable
model is kept distinct from the universal embedded curve.

% BEGIN PRESERVED PART 33-first-relations-v147.tex
\subsection{Universal horizontal specialization}
The new primary results are Theorems~\ref{thm:horizontal-universal-v166},
\ref{thm:all-arcs-v166}, \ref{thm:extension-class-v166}, and
\ref{thm:ramified-v166}. The complete companion papers supply their
focused reading order; this master retains every preceding proof.
